\subsection{Interior Milnor fibers and the surface inequality}
\label{sec:surface-estimates}

We need two signs in the nearby Euler comparison.  The local smoothing of an
interior singularity contributes at least the value~$1$ contributed by a smooth
point.  On the other hand, the complement of a reduced boundary in an extra
vertical surface has nonnegative Euler characteristic because its log canonical
divisor is nef.  We prove these statements separately.

We use ordinary Euler characteristic for compact local Milnor fibers.
All Betti numbers below have rational coefficients.

\begin{lemma}[The interior correction]
\label{lem:interior-correction}
Let $q:(\mathcal V,x)\to(\mathbb C,0)$ be a flat holomorphic germ whose
scheme central fiber $(S,x)$ is a normal surface with an isolated singularity
and whose sufficiently small nonzero fibers are smooth.  If $F_x$ is a
sufficiently small Milnor fiber of $q$, then
\[
  \chi(F_x)\geq 1.
\]
At a smooth point of a reduced Cartier central surface in a normal threefold,
the local nearby Euler characteristic is~$1$.
\end{lemma}

\begin{proof}
Greuel--Steenbrink's theorem \cite{GreuelSteenbrink1983} gives
$b_1(F_x)=0$ for a smoothing of a normal isolated singularity; no
Gorenstein hypothesis is required.  We use the formulation in
\cite[Theorem~2.18]{Greuel2019}.
Choose a Milnor representative with smooth transverse boundary.  Its interior
is a smooth Stein surface, being a closed analytic submanifold of a small
open ball slice.  The boundary collar makes this interior homotopy equivalent
to the compact Milnor representative.  By the Andreotti--Frankel
dimension bound \cite{AndreottiFrankel1959}, a complex Stein surface has
the homotopy type of a CW complex of real dimension at most two; see also
\cite[Theorem~1 and the following remark~(2)]{HillNacinovich1993}.
Consequently $b_j(F_x)=0$ for $j>2$.  The Milnor fiber is nonempty, so
\[
  \chi(F_x)=b_0(F_x)+b_2(F_x)\geq1.
\]

For the last assertion, let $t$ denote the pulled-back base coordinate.
At the point in question, the quotient of the threefold local ring by $(t)$
is regular of dimension two.  Lifting generators of its maximal ideal and
adjoining $t$ shows that the total local ring has embedding dimension at most
three, hence is regular.  The same dimension count shows that $t$ is a regular
coordinate.  Thus $q$ is smooth there and its small local fiber is contractible.
\end{proof}

\begin{lemma}[The surface inequality]
\label{lem:surface-bmy}
Let $S$ be a normal projective surface over $\mathbb C$, and let $B$ be a
reduced effective Weil divisor, possibly empty.  Suppose that $(S,B)$ is log
canonical and that the $\mathbb Q$-Cartier divisor $K_S+B$ is nef.  Set
$S^\circ=S\setminus\operatorname{Supp}B$.  Then
\[
  0\leq (K_S+B)^2
    \leq 3e_{\mathrm{orb}}(S,B)
    \leq 3\chi_c(S^\circ),
\]
where $e_{\mathrm{orb}}$ denotes Langer's orbifold Euler number.
\end{lemma}

\begin{proof}
Nefness on the projective surface implies pseudoeffectivity and nonnegative
self-intersection.  The latter can also be checked after pulling back to a
projective resolution, where the divisor is still nef and has the same square.
The pseudoeffective form of Langer's inequality gives
$(K_S+B)^2\leq3e_{\mathrm{orb}}(S,B)$
\cite[Corollary~5.2]{Langer2003}.

It remains to compare the two Euler numbers.  By Langer's definition,
smooth points of the boundary have weight zero.  At every singular point of
the pair lying on that boundary the local orbifold number is also zero, since the boundary
has a coefficient-one branch and the pair is log canonical
\cite[Theorem~7.6]{Langer2003}.  For
$Z=\operatorname{Sing}S\setminus\operatorname{Supp}B$, the weighted formula
therefore reads
\[
  e_{\mathrm{orb}}(S,B)
    =\chi_c(S^\circ)
      +\sum_{x\in Z}\bigl(e_{\mathrm{orb}}(x;S,0)-1\bigr).
\]
Every term in the sum is nonpositive by the local upper bound
\cite[Definition~3.4 and Corollary~7.8]{Langer2003}.  This proves the last
inequality, including all singular-point corrections.
\end{proof}

For application, write the reduced special fiber of the model as
\[
  Y_0=S_0+\sum_a S_a,
\]
where $S_0$ dominates the original special surface and each $S_a$ is an extra
component contracted to its marked point.  Let $B_S$ be the reduced boundary
on a component $S$ supplied by adjunction, and put
$S^\circ=S\setminus\operatorname{Supp}B_S$.
By Proposition~\ref{prop:log-model} and Lemma~\ref{lem:surface-boundary},
every $S_a$ is projective, $(S_a,B_{S_a})$ is dlt and hence log canonical, and its log
canonical divisor is nef: all its curves are contracted by the morphism over
which the threefold log divisor is nef.  Lemma~\ref{lem:surface-bmy} consequently
applies to each extra component.  It need not apply to $S_0$.

Proposition~\ref{prop:full-exceptional-boundary} identifies the nearby
complements $Y_c\setminus H$ with the smooth punctured fibers
$X_c\setminus\Gamma$ for $c\ne0$.  Thus an interior singular point
of any $S_a$ is smoothed by the family.  Indeed the central fiber is locally that component alone, is
normal and reduced Cartier, and the nearby fibers there are smooth.  Its
contribution in Lemma~\ref{lem:interior-correction} exceeds the smooth-point
value by the nonnegative number $\chi(F_x)-1$.

Lemma~\ref{lem:boundary-nearby-zero} supplies the vanishing of nearby
Euler values at intersections of boundary components, and
Proposition~\ref{prop:nearby-euler-comparison} combines it with these
interior corrections.  Theorem~\ref{thm:beta-constant} will use that
comparison and the two signs above to force every extra logarithmic square
to vanish.  The orbifold boundary calculation in the surface inequality
concerns a different invariant and does not supply the nearby-Euler
vanishing.
